\chapter{Modeling The World}

In parallel with the initial growth of the Computation Center, throughout the sixties, a series of changes took place nearby on campus, a few hundred yards to the north-east at the corner of Speedway and 24th. In 1960 the new Engineering Laboratories Building, later WRW, was occupied by a variety of applied research groups. In 1968 they were formally organized as the Department of Aerospace Engineering and Engineering Mechanics. This provided the institutional framework for a new era of computational engineering and science, associated with the pioneering work of researchers like Tinsley Oden and Byron Tapley. Under their leadership the transition to sophisticated, large-scale systems modeling accelerated. Oden’s foundational contributions to the Finite Element Method provided the mathematical rigor necessary to simulate complex solid and fluid mechanics, while Tapley’s expertise in orbital mechanics and satellite geodesy pushed the boundaries of high-precision gravity mapping and sea-surface height analysis. 

During the seventies, modeling and simulation grew beyond the more limited point-solutions of the fifties and sixties. Instead of focusing on where oil might be based on a few drill sites, engineers began using Finite Difference and Finite Element Methods to simulate fluid flow through porous rock across an entire field. And this era saw the dawn of satellite geodesy, with models of the Earth’s geoid requiring large-scale matrix computations that weren't practical a decade prior. The theory of Plate Tectonics was still relatively young in the late sixties, and computational modeling allowed scientists to simulate the thermal convection in the mantle that drives crustal movement. 

These types of models achieved new levels of predictive power and were a step-change up in data and numerical throughput. By creating a digital twin of a physical system, whether it was the Earth's crust or a jet engine, scientists could run what-if scenarios. Tinsley Oden played a leading role in these events and advocated for the view that computational engineering and science were in fact a new Third Pillar of Science, alongside the older pillars of experimentation and theory, stating that “computational science lies at the intersection of computer science, mathematics, science, medicine, and engineering…extending the scientific method, and represents the single most important scientific advance in human history, forever transforming the way science discoveries are made and how engineering and medicine are done.” and “In the future, there will be applications in areas we can’t even dream about. This won’t be solved without computational sciences and modeling.” 

A unifying theme of Texas and UT Austin history has been the land itself, working the land and defending it, and everything to do with the geology, natural resources, and physical history of Texas. In the Space Age, the technologies around these pursuits grew rapidly and extended outward to Earth orbit, the Moon, and beyond. Computing power, artificial satellites, laser ranging, digital communications, and dedicated engineers and scientists connected together a surprisingly wide-spread network of computational research projects, covering new ways of modeling both the Earth’s subsurface geological features and their measurable effects extending well into outer-space. 
These research directions grew together, side-by-side with the computers, satellites, and ground-stations needed to gather the resulting flood of data and rapidly process very large mathematical models, all in an operational manner, deeply woven into the national high-tech infrastructure. By following the development of these intertwined strands, from the sixties on through the seventies, a Texas-specific overview of computational engineering and science is possible. In particular, the strategic evolution of satellite gravity modeling provides a little-known and fascinating portal into revolutionary advances, transitioning from Cold War imperatives to become the foundations of modern remote sensing and resource exploration. 

Intellectual pioneers in Texas such as Byron Tapley, Bob Schutz, Ray Duncombe, and Mary Wheeler soon began recognizing that what was needed was an exhaustive understanding of the Earth’s gravitational field and that satellite orbits are directly linked to modeling of the Earth’s gravity, geology, sea-surface, subsurface, and interior. Their research brought the subsurface into contact with satellites, forming a synthesis of computational fluid dynamics and geodesy. This intersection was driven by the need to bridge localized subsurface processes such as flows in reservoirs with macro-scale observations of mass transport derived from satellite missions. Satellite gravity missions measure temporal variations in Earth’s gravitational field, which correspond directly to changes in mass distribution, for example due to fluid movement. 
Many of these developments originated in the early sixties, driven by the US Navy’s need for high-precision navigational updates for its submarine fleet. To ensure the accuracy of inertial guidance systems onboard submarines, the Transit satellite system was deployed, utilizing doppler frequency measurements to provide all-weather positioning. The first Transit satellite was launched into orbit on 13 April 1960. Among other information, it provided confirmation of the Earth's asymmetrical shape and highlighted the inadequacies of contemporary knowledge of the Earth's gravitational field for the prediction of satellite orbits, and other important near-Earth ballistic trajectories. Such prediction was essential for Transit's navigational role, in which receivers would determine their own position by monitoring the doppler shift from a satellite of known orbit. 

Already in May 1961 it was noted that “Meeting the ultimate program goals for Transit thus requires considerable improvement in the present knowledge of these factors (roughly the shape and mass distribution of the earth). This is the primary remaining development challenge of the Transit program.” Navigators have always relied on the stars for the purposes of terrestrial positioning. Artificial satellites revolutionized this paradigm by providing radio signals, available day and night and regardless of weather conditions. However, this shift introduced a significant challenge. Unlike the natural stars, these artificial stars are moving rapidly and are subject to complex orbital perturbations. 
A satellite is essentially a body in free fall influenced by forces such as gravity anomalies, atmospheric drag, and solar radiation pressure. Variations in the Earth’s gravitational field mean that the satellite’s trajectory is a direct reflection of the planet's irregular mass distribution. The concept of using satellites as test particles for probing the Earth’s gravitational field quickly developed and UT Austin was on the frontier of this research. The convergence of astrodynamics and geophysics established a computational grand challenge, where the requirements of gravity modeling and subsurface modeling drove a new scale of computing problems. Through the lens of computational geodesy, the satellite serves as a remote sensor for the Earth's interior, linking orbital mechanics to subsurface exploration and the management of global resources. 

\section{TRANET And Ray Duncombe}

Even before Transit, the roots of satellite gravity modeling trace directly back to the launch of Sputnik and the catalyzing effects it had for new methods and new concepts. During an intense week immediately following its launch on October 4, 1957, Ray Duncombe and Paul Herget of the US Naval Observatory worked around the clock to calculate the satellite's precise orbit. Duncombe would later bring his Naval Observatory experiences to a long research career at UT Austin. The natural, historical association of navigation with maritime matters made the navy and astronomers at the Naval Observatory key centers of activity in the early days. And in many ways the same can be said for astronomy itself, as stellar navigation and time keeping in general are essentially practical astronomy, and all of these are traditional maritime matters. There are deep historical ties between astronomy and the sea, with historic observatories associated with major ports and navies, and Ray Duncombe represented just such a link between UT’s McDonald Observatory in West Texas and the Naval Observatory.

By observing the Sputnik radio signal and its shift in frequency as the satellite approached and receded from a fixed station, they were able to determine Sputnik’s orbital elements. The early days immediately after the launch were a frantic effort to bridge the gap between the military's tracking capabilities and the unexpected reality of the launch, and there was anxiety stemming from several critical technical and political factors. The tracking network, known as Minitrack, was designed to track satellites transmitting at 108 MHz, the standard established for the 1957 International Geophysical Year. Sputnik transmitted on 20 MHz and 40 MHz frequencies that Minitrack stations were not originally equipped to handle. The lower frequencies used by Sputnik were subject to significant refraction as they passed through the Earth's ionosphere. This created complex errors in the tracking data that Duncombe, an expert in astronomical corrections, had to account for. 

Minitrack, short for Minimum Trackable Satellite, is an interesting story in itself, as it foreshadowed the style of federal Big Science that UT Austin would become associated with, such as the Transit and gps systems. It was the first US satellite tracking network, originally developed by the Naval Research Laboratory to track the Vanguard satellite. Minitrack stations were ground-based facilities designed to create a north-south fence of radio detection that a satellite must inevitably cross as it circled the Earth. Because it was designed solely to track satellites broadcasting at specific frequencies, the system had to be radically and rapidly altered to track Sputnik. 

The primary stations were strung along the 75th meridian, extending from Blossom Point Maryland down to Santiago Chile. This arrangement ensured that a satellite in an equatorial or semi-equatorial orbit would be observed by at least one station on each revolution. Unlike radar, which bounces a signal off an object, Minitrack was a passive system that listened for a radio signal emitted by the satellite. It used radio interferometry, a technique that measured the phase difference of radio signals arriving at pairs of antennas spaced specific distances apart. By comparing when the signal wave reached one antenna versus another, the system could determine the precise angle toward the satellite.
 
To track Sputnik, the Minitrack network had to undergo a frantic, improvised conversion. Technicians and station crews worked around the clock to modify their equipment to receive the lower Soviet frequencies. While experts predicted this conversion would take months, station crews accomplished it in roughly two weeks. At the Lima Peru station, for example, the crew hastily constructed 40 MHz dipole antennas using "2x4 end supports and lacing cord" and mounted them directly onto the existing antenna arrays. While Minitrack station crews were jury-rigging antennas, Duncombe was improvising mathematically. Instead of relying solely on nascent electronic computer processing, he performed a graphical solution of the satellite's orbit entirely by hand on a drawing table. This manual determination provided the first American calculation of the Soviet satellite's path. 

Meanwhile Duncombe was also involved in the numerical modeling and computing efforts needed to generate precise results. Because the Minitrack interferometric angle measurement was compromised by the frequency mismatch and ionospheric refraction, the work relied heavily on doppler shift measurements. By analyzing how the frequency of Sputnik's signal changed as it approached and receded, researchers could calculate the satellite's speed and closest approach, allowing them to mathematically reconstruct its orbit. This emergency adaptation of Minitrack, combined with data from visual observers and amateur radio operators, allowed the US to generate a definitive orbit for Sputnik within days of its launch. 

The calculations were performed at the Vanguard Computing Center in Washington using an IBM 704 mainframe. This vacuum-tube behemoth required a constant stream of data to refine the orbit. The team worked for 154 hours of continuous operation, roughly 6.5 days, performing over six billion calculations to generate a definitive ephemeris. Instead of using a known orbit to find a ship's position, as in the Transit system Duncombe later worked on, they had to use the signals received at known ground locations to deduce the unknown path of the satellite. The process was chaotic. The computing center was flooded with data from diverse and often unreliable sources, including visual observers, amateur radio operators, and experimental direction-finding stations. The team had to manually transcribe this data onto punched cards for the IBM 704 while managing sign errors and transmission discrepancies that could throw off the predicted orbit significantly. 

Another similar story unfolded a few months later, in 1958, with the second launch of an American satellite. There’s an uncanny resemblance and spirit of the times captured in this Richard Feynman legend.

\begin{quotation}
Explorer II, bearing a cosmic-ray detector that pushed its weight up to thirty-two pounds, soared skyward five weeks later and disappeared into the clouds. An army team watched under the guidance of Wernher von Braun, resilient veteran of the Nazi rocket program at Peenemünde. They listened to the fading rumble of the rocket and the rising beep of the radio signal transmitted to their squawk box. All seemed well. A half hour after the launch, they held a confident news briefing. Across the continent, where the Jet Propulsion Laboratory in Pasadena served as the army’s main collaborator in rocket research, a team was struggling with the task of tracking the satellite’s course. They used a room-size IBM 704 digital computer. It was temperamental. They entered the primitively sparse data available for tracking the metal can that the army’s rocket had hurled forward: the frequency of the radio signal, changing doppler-fashion as the velocity in the line of flight changed; the time of disappearance from the observers at Cape Canaveral; observations from other tracking stations. The JPL team had learned that small variations in the computer’s input caused enormous variations in its output. Albert Hibbs, the laboratory’s young research chief, had complained about this difficulty to his former Caltech thesis adviser: Feynman. Feynman bet that he could outcompute the computer, if fed the same data at the same rate. So when Explorer II lifted off the pad at 1:28 P.M., he sat in a JPL conference room, surrounded by staff members rapidly sorting the data for the computer. At one point Caltech’s president, Lee DuBridge, entered the room and was startled to see Feynman—who snapped, Go away, I’m busy. After a half hour Feynman rose to say he was finished: according to his calculations the rocket had plunged into the Atlantic Ocean. He left for a weekend in Las Vegas as the trackers kept trying to coax an unambiguous answer from their computer. Tracking stations at Antigua and Inyokern, California, persuaded themselves that they had picked an orbiting satellite out of the background noise, and “moonwatch” teams in Florida spent the night watching the skies. But Feynman was right. The army finally announced at 5 o’clock the next afternoon that Explorer II had failed to reach orbit.
\end{quotation}

Later, once the Transit system became operational in the early sixties, a network of dedicated doppler ground stations grew up, in parallel with the older Minitrack network of interferometric ground stations that Dumcombe had used for Sputnik orbit determination. This was the Transit network or Tranet, and UT Austin became deeply involved with operating two Tranet stations, one in Austin and the other at McMurdo Sound Antarctica. There was not a direct transition where Tranet entirely replaced Minitrack. They were two distinct global tracking networks that utilized different scientific principles, served different primary agencies, and were operated concurrently by different organizations, such as New Mexico State Universities Physical Science Laboratory, associated with the White Sands Missile Range and Fort Bliss in El Paso Texas. Minitrack was developed from 1956 for the Naval Research Laboratory and later transferred to NASA. It relied on angle tracking to determine a transmitting satellite's azimuth and elevation using massive, fixed antenna arrays deployed across fourteen global sites. PSL's Electromagnetics Section continued to support and calibrate the Minitrack network into the seventies, and it was ultimately replaced by modern tracking systems like NASA's Tracking Data and Relay Satellite System.

While Minitrack was originally designed to track early satellites like Vanguard, researchers realized that measuring a satellite's doppler shift could provide extremely precise orbital parameters. In late 1958, experimental doppler tracking loop of the Navy’s Vanguard I satellite proved that the doppler principle could achieve the sub-meter precision required for the Transit system to function, laying the foundation for Tranet. As satellite tracking matured, the doppler method proved highly effective for geodetic surveying. During the Southeastern United States Survey, the Navy's doppler method was evaluated against other basic satellite surveying methods, such as the SECOR system and cameras, and was determined to be the most accurate and reliable.

At the same time, throughout the seventies, UT Austin was creating and operating a very different type of satellite tracking station in West Texas at McDonald Observatory, utilizing the sophisticated new technique of satellite laser ranging or SLR. UT’s Tranet station at McMurdo and SLR station at McDonald put it at the frontier of gathering the new types of ground-based observational data feeding into the new computational models developing in Austin. The McMurdo and McDonald stations were extraordinary technical achievements and helped to establish a foundation of operational computing at UT that was strongly tied to the physical reality of satellites and cutting-edge radio and laser engineering. Both ground stations are well worth exploring in depth here, McMurdo Tranet first, followed by McDonald SLR, to make concrete the nature of these data streams at their sources.

It’s worth noting the curious blending of navigation, maritime, geology, astronomy, aerospace, and electrical engineering surrounding McMurdo and McDonald. In practical terms on the UT campus, for administrative reasons, Tranet was more associated with engineering and SLR was more associated with astronomy. Both influences blended together over time as part of what became the Center for Space Research, and this was one of CSR’s unique characteristics, with researchers from various departments collaborating freely. Researchers at CSR and ARL worked to integrate the Tranet doppler measurements with the SLR data collected at McDonald. 

By performing rigorous Precision Orbit Determination, CSR successfully resolved scale discrepancies, utilizing the doppler-derived geocentric coordinates to strengthen the highly accurate optical SLR reference frame. It was precisely this free-flowing, cross-departmental collaboration that allowed UT Austin to synthesize maritime-focused tracking with astronomical observations, establishing foundational frameworks for the modern International Terrestrial Reference Frame, International Celestial Reference Frame, and the International Earth Rotation Service to tie the ITRF and ICRF together.
Looking back, this was all a direct continuation of the long historical association between maritime navigation and astronomy, as Tranet was associated with the navy, while SLR was associated with the telescopes and dark skies of McDonald Observatory, and both were deeply intertwined with geodesy and geology. The historical ties to maritime navigation and geology were deeply interwoven into these tracking stations. The entire Transit Tranet system was originally funded and developed to solve a maritime problem: providing high-precision, all-weather inertial guidance updates for the U.S. Navy’s Polaris ballistic missile submarine fleet. The orbital models and benchmark coordinates generated by this military network ultimately laid the foundation for modern geological and cryospheric research. For example, the UT Institute for Geophysics later relied on the geodetic precision pioneered by these networks to conduct geophysical flights out of McMurdo, using airborne sensors and continuous GPS tracking to map bedrock topography and sub-ice marine environments across Antarctica.

The blending of astronomy and aerospace was perfectly embodied by Ray Duncombe, whose career bridged the gap between optical telescopes and satellite tracking. Duncombe studied astronomy and spent decades at the Naval Observatory working with classical optical transit circle telescopes to map stellar positions. He then applied this classical celestial mechanics expertise to determine the orbits of early artificial satellites like Sputnik and Vanguard. His career spanned an era covering the dawn of electronic digital computing and the ensuing achievement of unprecedented levels of precision, eventually allowing measurements with centimeter-level accuracy for tracking tectonic plate motion, Earth's polar motion, and variations in the gravitational field.

When Duncombe joined UT Austin in 1975, it was as a professor in the Department of Aerospace Engineering and Engineering Mechanics, and as a member of CSR. Characteristically, he physically connected the engineering curriculum with astronomy, geodesy, and navigation. He regularly took the aerospace graduate students from his Determination of Time course on field trips to McDonald Observatory. Since this was an eight-hour drive from Austin and required multiple days, the seriousness of the topics involved is obvious. West Texas and McDonald were special places for learning about the realities of celestial time keeping and navigation, with some of the darkest skies easily available, though Antarctica was of course in a very different and much more difficult and isolated category.

During the seventies, the Transit satellite tracking station at McMurdo Sound Antarctica, designated as Station 019, operated as a critical polar node within the Navy's global Tranet system. Because polar orbits are fundamental to navigation constellations, high-latitude observations from McMurdo were essential for accurately determining satellite inclination and nodal progression. A distinctive hallmark of this station was its unconventional staffing paradigm. Rather than relying on military personnel or career technicians, the station was operated by a two-person team comprising one graduate and one undergraduate student from UT Austin. Selected primarily from electrical engineering, these students brought the expertise in radio frequency electronics, digital logic, and communication systems necessary to manage the remote outpost. Maintaining a continuous, year-round presence in Antarctica required complex institutional coordination among UT, the Navy, and the National Science Foundation.

Tracking operations at McMurdo began on February 5, 1965, within a National Science Foundation building, and were originally managed by New Mexico State University's Physical Science Laboratory. After PSL withdrew in December 1966 due to administrative issues, UT Austin's Applied Research Laboratories officially reopened the station on October 10, 1968. ARL maintained continuous management of the site for a quarter-century. Aside from satellite tracking, McMurdo Sound was a hub for various other scientific projects under the U.S. Antarctic Research Program. During the 1975 to 1976 season, researchers used McMurdo tracking facilities as a master translocator for roving satellite receivers to map ice sheet movements. Biologists and geologists also studied the area extensively, investigating topics such as the ocean-floor communities of McMurdo Sound, the paleoclimatic history of the region via borehole geology, and the unique "antifreezes" that prevent native Antarctic fishes from freezing in the ice-laden waters.

Operating the McMurdo Tranet station presented severe technical and environmental challenges that demanded a high degree of innovation from its student staff. The students assumed responsibility for the entirety of the station's technical ecosystem, which included operating and maintaining complex doppler receivers and antennas to maximize data yield from every satellite pass. Furthermore, they were tasked with vital infrastructure maintenance, ensuring that specialized equipment and electrical generators continued to function in an extreme environment where temperatures frequently plummeted below -50 degrees Celsius. In early 1971, the station's tracking operations were automated, allowing for more consistent and higher-density data collection without requiring constant human intervention at the consoles. However, the students were still responsible for operating and maintaining the complex doppler receivers and antennas to ensure maximum data yield from every satellite pass. 

A massive logistical hurdle was getting the collected orbital data from the Antarctic to analysts at Dahlgren, Virginia. Throughout the 1970s, the station relied on high-frequency HF radio links to relay data to New Zealand, which then forwarded the packets to the United States via AUTODIN Automatic Digital Network circuits. Because HF communications in Antarctica were notoriously unreliable due to "polar blackout" events caused by solar flares, the students had to become experts in radio propagation, carefully selecting optimal frequencies and times for data transmission.
Beyond its operational role as a passive data collector, the McMurdo facility functioned as an active research site for ionospheric propagation studies sponsored by the Navy and the Defense Mapping Agency. Because the Earth's magnetic field lines converge at the poles, the polar ionosphere experiences intense disturbances from cosmic radiation and solar wind, causing significant refraction of satellite signals. UT Austin students conducted year-round investigations into these phenomena, utilizing the stable signals of Transit satellites to map variations in ionospheric density. Their findings were instrumental in developing the dual-frequency correction methods, transmitting simultaneously on 150 MHz and 400 MHz, that allow receivers to calculate and cancel out ionospheric delays. Ultimately, the high-precision doppler measurements and original research conducted by these students in one of the world's most hostile environments established the geodetic benchmarks that refined the World Geodetic System, WGS 84, laying the foundational coordinates for the modern Global Positioning System.

\section{McDonald Observatory And West Texas}

Laser ranging played a core role in the birth of modern space geodesy, and during the seventies the McDonald Observatory in West Texas established a virtual hegemony over these types of observations, beginning with the transformation of the Moon from a subject of celestial observations into a high-precision geodetic reference target. In space geodesy, the Earth is understood as a highly dynamic system rather than a rigid sphere. Calculating the precise round-trip travel time of a laser pulse to the Moon or an artificial satellite requires an extraordinarily accurate mathematical model. To isolate the distance measurement, researchers must account for an array of Earth's dynamical influences such as overall rotation rate, precession, nutation, and polar motion. With the placement of the first laser retro-reflector arrays on the lunar surface, particularly during the Apollo program, researchers were able to revolutionize lunar distance measurement. The Apollo era retro-reflectors enabled centimeter-level measurements from McDonald, replacing radar uncertainties of 0.7 miles during the sixties, and early laser experiments from 1965 that achieved only 180-meter accuracy.

The Moon was a difficult laser ranging target because of its distance and the orbital velocity of the Moon’s relative to the Earth. The phenomenon of velocity aberration dictated that smaller retro-reflector cubes would provide higher intensity for a single ground-based receiver, so that arrays of smaller cubes were preferable. These physical constraints necessitated a permanent tracking station capable of capturing the few photons available after the enormous losses compared to retro-reflectors on near-Earth satellites. The small signal called for a large telescope, and this was one of the key motivations for NASA to fund construction of the new 107-inch telescope at McDonald in support of the Apollo program. The 107-inch telescope served as the world’s only routine lunar ranging facility for over fifteen years, achieving its first successful returns on August 19, 1969. 

During the seventies, satellite laser ranging at the McDonald Observatory developed in parallel with the lunar laser ranging program. While the 107-inch telescope is most famous for its post-Apollo lunar tracking, its observing time was also shared with artificial-satellite ranging. UT Austin's Center for Space Research utilized the SLR data collected to make pioneering advances in precision orbit determination and geophysical measurement. Key satellite missions and SLR milestones included the Geodynamics Experimental Ocean Satellite, which was used to study precision orbit determination and to develop regional gravity models. With Lageos in 1976, a high-density passive reflector in medium Earth orbit, CSR researchers were among the first to demonstrate that SLR could be used to measure tectonic plate motion. SLR tracking of Lageos also allowed CSR to provide definitive results for Earth's polar motion, calculate the value of Earth's mass, and make the first measurements of seasonally-induced time-variable gravity. And with SeaSat in 1978, one of the earliest active Earth-observing satellites, CSR led the precision orbit determination team tasked with using tracking data to support radar altimetry processing for accurate ocean surface topography and gravity measurements. 

As the demand for artificial satellite tracking grew, the engineering and operational footprint at McDonald Observatory expanded. In the late seventies, UT Austin developed and tested the first Transportable Laser Ranging System at the McDonald site, validating coordinate determination algorithms and demonstrating the viability of mobile laser tracking. This effort evolved into the dedicated McDonald Laser Ranging Station sited near the 107-inch, while starting in 1979, various NASA-supplied satellite ranging stations began operating continuously at the observatory. The connection between the Transit satellite system and McDonald Observatory as a geodetic station lies in the computational modeling at UT’s CSR and ARL. These models combined Transit's radio-frequency doppler tracking with the optical SLR networks anchored at McDonald, creating an early end-to-end tie between the celestial frame, terrestrial frame, and the Earth’s position and rotation.

The Transit satellites continuously broadcasted two coherent carrier frequencies at 150 MHz and 400 MHz to allow ground receivers to eliminate first-order ionospheric refraction. The Tranet global tracking network supporting these geodetic investigations consisted of stations equipped with ultrastable rubidium or cesium frequency standards. UT Austin served as a vital institutional node linking this doppler-based satellite network with high-precision laser ranging. ARL developed specialized mathematical packages to compute ionospheric and tropospheric refraction corrections for Tranet. To standardize geodetic data archiving and coordinate exchange, ARL developed the FICA Floating Integer Character ASCII format, which served as the standard for geodetic GPS and Transit data before the universal adoption of RINEX.
The direct operational connection between the Transit system and McDonald Observatory was forged through modeling and orbital synthesis at CSR. Precision orbit determination for major oceanographic altimetry missions required integrating data from the 48-station Tranet network with optical tracking data. CSR researchers merged the Tranet doppler measurements with SLR data collected by McDonald and similar observatories. To resolve scale discrepancies, CSR executed precise coordinate ties between the Tranet receivers and the SLR network. This synthesis successfully tied doppler-determined geocentric coordinates to the highly accurate SLR reference system, defined by targets such as Lageos, reducing radial orbit errors to approximately 20 cm. ARL and CSR also evaluated prototype GPS geodetic receivers by comparing GPS pseudorange and phase observations against legacy Transit integrated doppler phase accumulations. These field tests demonstrated sub-meter baseline determination, validating the transition from Transit-based datums to the modern WGS 84 World Geodetic System and the ITRF International Terrestrial Reference Frame, a framework to which McDonald Observatory continuously contributes.

As McDonald developed into a key geodetic reference site across the seventies, the nature of its West Texas geology was crucial. For one thing, it is onboard one of the major cratons that have played major roles across geologic deep time, though near a margin or boundary zone rather than in an ultra-stable core region. In other words, it is in a fairly stable and ancient continental plate region. And as importantly, it is sitting on deeply stable rock, at the opposite end of the spectrum from a site on top of subsurface reservoirs which can expand and contract. 

The Davis Mountains represent the largest contiguous volcanic remnant of the mid-Tertiary Trans-Pecos Magmatic Province, a massive alkalic volcanic field formed over forty million years of tectonic reorganization. Between 39 and 35 million years ago, a violent series of caldera collapses and explosive eruptions buried the region in vast sheets of ash-flow tuffs and flood rhyolites. The resulting stratigraphic architecture produced a landscape characterized by low-viscosity, high-temperature silicic magmas that cooled into exceptionally dense, structurally rigid igneous rocks. Today, these precise geomechanical properties serve as the foundational bedrock for the world-class optical and geodetic instrumentation at the McDonald Observatory.

The volcanic rocks supporting Mount Locke and Mount Fowlkes are highly fractured. Rainwater rapidly seeps deep into these fractures, accumulating as "rock moisture" within the unsaturated zone. During heavy monsoon seasons, the increased water mass within the fractured trachyte and tuff increases the local gravitational pull. This localized subsurface water storage mimics the gravimetric signature of actual tectonic crustal subsidence. The deep infiltration of water into the observatory's bedrock plays a critical role in the broader regional hydrogeology of Far West Texas. The precipitation that falls on the Davis Mountains percolates underground, migrating through highly porous volcanic pathways and underlying Cretaceous limestone layers to recharge artesian desert spring systems, most notably the San Solomon Springs in Balmorhea. While the continuous baseflow of these springs originates distally from the Salt Basin Bolsons to the west, the system is augmented multifold by local stormflow recharge moving through the subsurface from the Davis Mountains. The fractured volcanic layers, particularly the basal conglomerate of the Huelster Formation and associated vesicular lavas, act as rapid transport conduits, though some volcanic layers may also retain water and release it slowly over months, sustaining the springs long after precipitation events.

While the interior of the Davis Mountains is relatively stable, the region is bounded by major active fault zones associated with the Rio Grande Rift and the Texas Lineament, a major crustal boundary separating the stable North American craton from the active Chihuahuan borderlands. The persistent seismic risk of the subsurface was demonstrated by the 1931 Valentine Earthquake, the largest in Texas history, centered just 25 miles west-southwest of the observatory. Although the dense, consolidated volcanic sequences of Mount Locke and Mount Fowlkes act as a rigid, dampening block that mitigates catastrophic ground failure, the continuous, low-frequency seismic stress of the basin-and-range faults must be continuously accounted for in the observatory's geodetic and astronomical measurements. McDonald Observatory’s enduring scientific legacy is inextricably linked to its geology and subsurface environment. In recognition of this intrinsic connection, its mission has evolved to include geological, hydrological, and ecological research across its 650 acres of Chihuahuan Desert terrain.

The telescope site’s geomechanical stability was paramount. High-precision astronomical instrumentation requires structural foundations entirely decoupled from wind shear, thermal expansion, and human-induced vibrations. The Mount Locke Formation, restricted to a 93-meter-thick flow of coarse, highly porphyritic metaluminous trachyte, provided an exceptionally dense and unyielding bedrock. By anchoring the 107-inch telescope’s massive 160-ton structure directly into this rigid volcanic substrate, engineers achieved the absolute dimensional stability necessary to isolate the optical path from environmental noise.

The rigid bedrock of the Mount Locke Formation provided the exceptional load-bearing strength necessary for the telescope's foundations and allowed for the site’s calibration to the International Terrestrial Reference Frame and the International Celestial Reference Frame with zero margin for error. The seventies data boom allowed for the first empirical measurements of Earth Orientation Parameters and the determination that the North American Plate drifts southwestward across the mantle.

The data boom resulting from the network of lunar reflectors and the unyielding trachyte foundation of Mount Locke allowed McDonald’s LLR measurements to serve as a highly stable anchor for the ITRF, ICRF, and IERS. Through continuous tracking, researchers successfully utilized the LLR data to empirically measure Earth Orientation Parameters, including polar motion, nutation, and variations in the length of day. By definitively linking terrestrial coordinates to the Earth's center of mass, the observatory tracked the continuous drift of the North American Plate across the mantle, validating the mechanics of plate tectonics and cementing the site's legacy in the geosciences. McDonald Observatory and the Center for Space Research were the first laser ranging group to provide operational observations of these Earth Orientation Parameters, notably polar motion and length of day. Through the accumulation of SLR and LLR data, researchers have been able to isolate variations in the Earth's principal figure axis and key geopotential coefficients across timescales ranging from sub-daily to decadal.

Along with SLR, researchers tracked the precise distances between different global ground stations. Over time, these inter-station baseline rates reveal the slow, continuous drift of the Earth's tectonic plates with a precision of better than one centimeter. McDonald Observatory played a critical role in this endeavor and accumulated a continuous legacy of providing reference positions that accurately tracked the motion of the North American Plate from the seventies onwards. These measurements are a cornerstone of the ITRF for tracking the dynamic movement of the Earth's crust.

Understanding continental motion is an absolute prerequisite for accurately measuring global sea levels. CSR led the Precision Orbit Determination efforts for satellites monitoring variations in sea surface height using radar altimeters such as the TOPEX/Poseidon and Jason series missions. These altimeters measured the distance between the satellite and the ocean surface. To determine the actual height of the sea, CSR tracked the satellite's exact altitude and position in space relative to the center of the Earth. This required tracking the satellite from ground stations. Because these ground stations are constantly moving due to tectonic drift, their exact coordinates must be continuously updated using the ITRF. If continental motion were not accounted for, the changing position of the ground stations would introduce significant errors into the satellite's orbit calculations, which would in turn corrupt the sea level measurements. While altimeter missions track the total height and volume of the continental and ocean surfaces, CSR has also led efforts to directly track subsurface and oceanic mass variations.

Observations from McDonald Observatory and the global geodetic network reveal that Earth's rotation is altered by three primary mechanisms: external gravitational torques, internal mass redistribution, and the transfer of angular momentum between the solid Earth and its fluid layers. The gravitational attraction of the Sun and Moon generate oceanic and solid-body tides, and the friction from these tides causes the Earth's rotation rate to slowly decrease. To conserve the total angular momentum of the Earth-Moon system, this deceleration forces the Moon to spiral away from the Earth at a rate of approximately 3.8 centimeters per year.

As the Earth displaces its mass, its moment of inertia changes, subsequently altering its spin rate and the position of its poles. Variations in the Earth's length of day on decadal scales are heavily influenced by the internal structure of the planet. Gravitational and hydrodynamic coupling between the fluid outer core, the solid inner core, and the mantle exerts torques at the boundaries of these layers, leading to persistent oscillations in Earth's rotation rate. Space geodetic observations of polar motion revealed that around 2005, the Earth's pole began drifting abruptly toward the east. The precision of LLR is so extraordinary that even biospheric cycles are detectable. As tree sap rises or falls seasonally in the heavily forested Northern Hemisphere, the corresponding change in Earth's moment of inertia minutely alters the planet's rotation rate, a signature captured in the ranging data.

\section{Satellites To Subsurface}

There was an underlying relationship that birthed both satellite navigation and satellite geodesy. If a known ground position can resolve an unknown orbit, then a precisely determined orbit can resolve an unknown ground position. The central challenge of the era became the minimization of orbital errors, as any uncertainty in a satellite’s velocity or position translated directly into inaccuracies in a navigational fix. This requirement for precision tended to shift the focus from navigation to geodesy, where the satellite was reimagined as a test particle probing the Earth’s gravity field and helping to map gravitational irregularities. The successful use of artificial satellites as purely gravitational test particles hinged entirely on isolating them from non-gravitational perturbations. 

Because a satellite in orbit is constantly pushed and pulled by atmospheric drag, solar radiation pressure, and thermal radiation emitted by the Earth, engineers had to design spacecraft that would be as immune to these surface forces as physically possible. This requirement led directly to the cannonball satellite design, which prioritized an extremely low area-to-mass ratio. By making the satellites as small, spherical, and dense as possible, engineers minimized the surface area vulnerable to drag and solar pressure while maximizing the mass that the Earth's gravity could act upon. Launched by France in 1975, Starlette was a small sphere only 14 centimeters in diameter, but it contained a core of depleted uranium to artificially increase its mass. This gave it an incredibly small area-to-mass ratio, ensuring its trajectory was almost entirely dictated by gravity rather than drag. Commissioned by NASA, the Lageos Laser Geodynamics Satellite satellites were built as highly dense aluminum spheres surrounding heavy brass cores.

To help minimize the orbital errors and extract precise geodetic data, these test particle satellites were covered in optical retro-reflectors and tracked from the ground using satellite laser ranging. Unlike early radio or optical tracking, SLR provided an absolute, unambiguous measurement of the distance between the ground station and the satellite's center of mass, which was crucial for locating the orbit precisely relative to the Earth's geocenter. Eventually the success of these cannonball satellites shifted the central challenge of astrodynamics once again as the collected SLR data was used to create accurate maps of the Earth's gravity field. Because the gravitational errors had been so drastically reduced, subtle non-gravitational surface forces like the minute push of solar and Earth radiation pressure re-emerged as the primary limitation in achieving sub-centimeter orbital precision for modern satellite missions.

At UT, the UTOPIA software running on the Computation Center CDC 6600 was specifically designed to perform Precision Orbit Determination on these cannonball test particle satellites. By tracking exactly how these dense spheres fell through space, researchers were able to invert the data by numerically solving very large matrix systems and mapping the density variations beneath the Earth's crust. Researchers began to quantify the subtle variations in the planet's mass and to use the Earth's gravity field as a means for imaging its internal structures, much like a vast CAT or MRI scanning machine. Future evolutionary leaps in this technology involve using laser-cooled cesium atoms as the falling proof mass inside interferometeric sensors that will be a million times more sensitive than current hardware. This technology is expected to "essentially render the Earth's surface transparent" and produce incredibly detailed three-dimensional maps of underground structures and track the movement of magma before volcanic eruptions or pinpoint highly specific mineral and oil deposits.
The extraction of geodetic information from new technologies such as Tranet and satellite laser ranging necessitated a computational infrastructure capable of processing unprecedented data volumes. By the seventies, gravity modeling had emerged as a grand challenge for computational science, alongside subsurface modeling and weather modeling. UT Austin became a central hub for this research, with the Computation Center playing its role in the massive data problems inherent in satellite geodesy, and the early ties with the oil industry in Houston naturally carried forward into the new era. This connection was strengthened through Rice University in Houston which continued working with the Computation Center to facilitate subsurface reservoir modeling. The collaboration between Mary Wheeler at Rice and the computational efforts at UT Austin and Exxon highlighted the shared mathematical underpinnings of these fields, involving complex modeling and simulation.
